\section{总结与延伸}

\subsection{核心经验汇总}

OpenAI 团队的 Harness Engineering 实验揭示了 Agent-First 开发范式的若干核心原则：

\begin{table}[H]
\centering
\small
\begin{tabular}{p{4cm}p{9cm}}
\toprule
\textbf{原则} & \textbf{内涵} \\
\midrule
人类驾驭，Agent 执行 & 工程师从写代码转向设计环境、表达意图、构建反馈循环 \\
地图优于百科全书 & AGENTS.md 应是简短的目录，知识库分布式存储 \\
可见性决定存在性 & Agent 在运行时看不到的东西等于不存在 \\
约束是速度的前提 & 严格的架构约束使高速产出不致衰退 \\
修正廉价，等待昂贵 & 高吞吐量下快速合并、后续修正优于阻塞等待 \\
持续垃圾回收 & 技术债务小额持续偿还，不让其复利累积 \\
无聊技术更可靠 & 稳定、简单、训练集覆盖充分的技术对 Agent 更友好 \\
\bottomrule
\end{tabular}
\caption{Harness Engineering 核心原则总结}
\end{table}

\subsection{开放问题}

文章坦诚地指出了几个尚未回答的问题：

\begin{enumerate}[nosep]
  \item \textbf{长期架构一致性}：在完全 Agent 生成的系统中，架构的一致性如何随年数演进？五个月和五年是完全不同的时间尺度。
  \item \textbf{人类判断的最优杠杆点}：人类判断在哪些环节提供最大的杠杆效应？如何编码这些判断使其能够持续复利？
  \item \textbf{模型能力迭代的影响}：随着模型持续变得更强大，这套系统如何演进？当前的很多约束是否会变得不再必要？
\end{enumerate}

\subsection{对从业者的启示}

\begin{knowledgebox}{从这篇文章中可以立即行动的三件事}
\begin{enumerate}[nosep]
  \item \textbf{审视你的 AGENTS.md / CLAUDE.md}：它是地图还是百科全书？如果超过 200 行，考虑拆分为结构化的 docs/ 目录。
  \item \textbf{将隐性知识推入仓库}：那些只存在于 Slack、会议纪要或人们脑中的架构决策，值得被文档化为版本控制的 Markdown。
  \item \textbf{投资自定义 Linter}：特别是带有修复指导的错误消息，它们是``一次编码、处处生效''的乘数效应工具。
\end{enumerate}
\end{knowledgebox}

\subsection{延伸阅读}

\begin{itemize}[nosep]
  \item OpenAI Codex 官方文档：\url{https://openai.com/codex}
  \item Anthropic 关于 CLAUDE.md 最佳实践的讨论
  \item Martin Fowler 关于``Strangler Fig Pattern''的文章——渐进式重构与 Agent 垃圾回收有相似的哲学
  \item ``Choose Boring Technology'' (Dan McKinley, 2015)——文章中``无聊技术对 Agent 更友好''的论点的早期理论基础
\end{itemize}

\vspace{1em}

\begin{center}
\textit{``Building software still demands discipline, but the discipline shows up more in the scaffolding rather than the code.''}

--- Ryan Lopopolo, OpenAI
\end{center}

%% --- Fin del contenido --- %%

\end{document}
